\subsection{Sprint 4}

A Sprint 4 representou a etapa final de consolidação do projeto. As atividades previstas envolviam o card de projetos apoiados no dashboard da empresa, a criação da tela de notificações no painel dos administradores, a criação do perfil da empresa e a remoção dos mocks restantes do CompanyDashboard. Todas as atividades previstas foram concluídas, o que ajudou a aproximar o sistema de um estado mais consistente para apresentação e validação final.

Nesta sprint, minha atuação esteve bastante relacionada a revisão, planejamento e implementação de ajustes no painel administrativo. Participei da contabilização das horas da weekly de sábado, revisei pull requests e ajudei colegas com erros encontrados durante o desenvolvimento. Também trabalhei no planejamento e implementação de um débito técnico no painel dos administradores, buscando reduzir inconsistências e tornar essa área mais coerente com as funcionalidades reais do sistema.

Uma das principais frentes em que atuei foi a criação do planejamento de implementação da listagem de notificações no painel de admin. Para isso, realizei uma varredura do painel administrativo, listei inconsistências de design, participei da weekly e alinhei decisões com integrantes da AGES III e IV. Em seguida, contribui para a implementação da listagem de notificações, conectando essa necessidade ao fluxo de mediação e acompanhamento de vínculos dentro da plataforma.

A sprint também envolveu melhorias importantes no dashboard da empresa. A remoção de mocks restantes do CompanyDashboard foi relevante porque aproximou a aplicação de um comportamento real, baseado em dados do backend. O card de projetos apoiados também contribuiu para tornar a experiência da empresa mais útil, permitindo visualizar de forma mais direta os projetos vinculados e a divisão entre modalidades de investimento.

Apesar das entregas concluídas, o maior problema encontrado foi a falta de motivação e engajamento da equipe de desenvolvimento. Ao final do projeto, ficou perceptível que parte do time já demonstrava menor energia para continuar evoluindo a aplicação. Esse desgaste pode ser explicado por diversos fatores, como fim de semestre, acúmulo de demandas, complexidade do projeto e perda gradual de interesse pelo produto.

No meu caso, essa sprint também evidenciou uma redução da minha própria presença no projeto. Eu ainda contribuía por meio de revisão, planejamento, apoio e implementação pontual, mas já não tinha o mesmo nível de envolvimento das fases iniciais. Essa redução não ocorreu por uma dificuldade técnica específica, mas por uma combinação de menor interesse pessoal pelo business do projeto e pela sensação de que eu já não estava aprendendo tantas coisas novas quanto em experiências anteriores.

A principal lição aprendida foi que a escolha de projetos deve considerar também preferências pessoais e interesse pelo domínio de negócio. Durante muito tempo, eu avaliei projetos principalmente pela stack ou pelo desafio técnico. No entanto, a Sprint 4 deixou claro que o business precisa despertar algum nível de curiosidade nos desenvolvedores, principalmente quando o projeto entra em fases menos empolgantes, como ajustes, revisões e correções. Quando não existe conexão com o problema que está sendo resolvido, manter o engajamento se torna mais difícil.

Como próximos passos, o projeto ainda demandava ajustes no deploy da aplicação. A infraestrutura já possuía uma base organizada, com Docker, AWS e separação de ambientes, mas ainda havia necessidade de estabilizar e ajustar pontos finais para garantir uma entrega mais confiável. Para mim, essa sprint encerrou a AGES III com uma reflexão importante: maturidade técnica não depende apenas de saber implementar, mas também de entender quais tipos de problema realmente despertam interesse e sustentam participação ao longo do tempo.
